\begin{examplebox}
\textbf{例 1：已知焦半径求准距。}
$x^2/16-y^2/9=1$ 上 $PF_R=10$。
$a=4,c=5,e=5/4$，所以 $d(P,\ell_R)=10/(5/4)=\boxed8$。
回验：右支可取 $x=56/5$，此时 $y^2=1539/25$，
到焦点 $(5,0)$ 的距离平方为 $(31/5)^2+1539/25=100$，到准线 $x=16/5$ 的距离为 $8$。

\medskip
\textbf{例 2：由横坐标转化距离。}
$x^2/4-y^2/12=1$ 上 $x=5$。
$a=2,c=4,e=2$，右准线 $x=1$。
故准距 $=|5-1|=\boxed4$，焦半径 $=2\cdot4=\boxed8$。
回验：$y^2=63$，到 $(4,0)$ 的距离平方为 $1+63=64$。
\end{examplebox}
\clearpage
\begin{examplebox}
\textbf{例 3：已知右焦半径，不能擅自限定右支。}
在 $x^2/4-y^2/5=1$ 上，$PF_R=6$，求横坐标。
$a=2,c=3,e=3/2$，右准线 $x=4/3$。
\begin{center}\begin{tikzpicture}[font=\small,>=stealth]
\draw[black!30,rounded corners=2pt] (0,0) rectangle (14.4,8.2);
\node[anchor=west,font=\small\bfseries] at (.5,7.72) {例 3：到右焦点距离为 $6$，两支上都可能有点};
\begin{scope}[shift={(6,4)},scale=.4]
\draw[blue,very thick,domain=-6:6,samples=161,variable=\t] plot ({2*sqrt(1+\t*\t/5)},\t);
\draw[blue,very thick,domain=-6:6,samples=161,variable=\t] plot ({-2*sqrt(1+\t*\t/5)},\t);
\draw[orange!85!black,thick] (3,0) circle (6);
\draw[violet,thick] ({4/3},-6.3)--({4/3},6.3);
\node[fill=black,circle,inner sep=1.8pt] at (3,0) {};
\node[below right=7pt] at (3,0) {$F(3,0)$};
\node[fill=black,circle,inner sep=1.8pt] at (5.3333333333,-5.5277079839) {};
\node[fill=black,circle,inner sep=1.8pt] at (5.3333333333,5.5277079839) {};
\node[right=8pt] at (5.3333333333,5.5277079839) {$x=16/3$};
\node[fill=black,circle,inner sep=1.8pt] at (-2.6666666667,-1.9720265944) {};
\node[fill=black,circle,inner sep=1.8pt] at (-2.6666666667,1.9720265944) {};
\node[left=12pt] at (-2.6666666667,1.9720265944) {$x=-8/3$};
\node[below left] at ({4/3},0) {$\ell$};\end{scope}

\node[anchor=west] at (.7,.55) {圆 $PF=6$ 的四个交点分属两支，对应两个横坐标};
\end{tikzpicture}\end{center}

由 $6=\tfrac32|x-4/3|$，得 $x=4/3\pm4$，即
\[
\boxed{x=16/3\quad\text{或}\quad x=-8/3}.
\]
回验：前者 $y=\pm5\sqrt{11}/3$，后者 $y=\pm\sqrt{35}/3$，
代入双曲线均成立，且到 $(3,0)$ 的距离平方均为 $36$。
只有题目另加“$P$ 在右支”时，才仅保留 $16/3$。
\end{examplebox}
\clearpage
\clearpage
\begin{examplebox}

\medskip
\textbf{例 1 补充：为什么不能由 $PF_R=10$ 就默认 $P$ 在右支？}
用 $PF_R=|\frac54x-4|=10$ 直接解得
\[
\frac54x-4=10\quad\text{或}\quad\frac54x-4=-10,
\qquad \boxed{x=\frac{56}{5}\ \text{或}\ x=-\frac{24}{5}}.
\]
两值分别满足 $x\ge4$ 和 $x\le-4$，且都能在双曲线上找到点。
左支取 $x=-24/5$ 时 $y^2=99/25$，
到 $(5,0)$ 的距离平方为 $(-49/5)^2+99/25=100$。
原题只问准距，所以两个分支\textbf{都得到 $8$}，不会影响原答案。

\end{examplebox}
\clearpage
\begin{examplebox}
\medskip
\textbf{例 4（新）：忘了准线公式，怎样从顶点现场找回？}
已知双曲线 $\dfrac{x^2}{9}-\dfrac{y^2}{16}=1$，求右准线。
由 $a=3,b=4,c=5,e=5/3$，右顶点为 $A(3,0)$，右焦点为 $F_R(5,0)$。
设准线为 $x=t$（$0<t<3$），顶点处
\[
\underbrace{AF_R}_{2}=\frac53\underbrace{(3-t)}_{\text{顶点到准线距离}}
\quad\Longrightarrow\quad
\boxed{t=\frac95}.
\]
\textbf{考场意义：} 无须死背 $a^2/c$，只要记得焦准比，代入顶点就可恢复。

\medskip
\textbf{例 5（新）：已知点到右焦点距离，求横坐标并筛根。}
在 $\dfrac{x^2}{9}-\dfrac{y^2}{16}=1$ 上，已知 $PF_R=7$。
由 $e=5/3$，$PF_R=|5x/3-3|$，所以
\[
\left|\frac{5x}{3}-3\right|=7
\Longleftrightarrow \frac{5x}{3}-3=\pm7
\Longleftrightarrow \boxed{x=6\ \text{或}\ x=-\frac{12}{5}}.
\]
但双曲线要求 $|x|\ge a=3$，于是 $x=-12/5$ \textbf{必须舍去}。
最终 $\boxed{x=6}$。回验 $y^2=16(36/9-1)=48$，
到 $(5,0)$ 的距离 $\sqrt{(6-5)^2+48}=7$。
\textbf{经验：} 绝对值开出两个根不代表两个根都在双曲线上，
必须检查 $|x|\ge a$。

\medskip
\end{examplebox}
\clearpage
\begin{examplebox}
\textbf{例 6（新）：从“焦点 + 准线 + $e$”反向求方程。}
固定 $F(2,0)$、$\ell:x=0$，求满足 $PF=2d(P,\ell)$ 的点的轨迹。
令 $P(x,y)$，由两边非负可平方：
\[
(x-2)^2+y^2=4x^2
\Longleftrightarrow 3x^2+4x-y^2=4
\Longleftrightarrow 3\left(x+\frac23\right)^2-y^2=\frac{16}{3}.
\]
整理为
\[
\boxed{\frac{(x+2/3)^2}{16/9}-\frac{y^2}{16/3}=1}.
\]
中心 $(-2/3,0)$，$a=4/3$、$b=4/\sqrt3$，
$c=8/3$、$c/a=2$；右焦点 $-2/3+8/3=2$，
右准线 $-2/3+(16/9)/(8/3)=0$，与条件吻合。

\medskip
\end{examplebox}
\clearpage
\begin{examplebox}
\textbf{例 7（新）：纵向焦准比的直接使用。}
若 $\dfrac{y^2}{9}-\dfrac{x^2}{16}=1$，
则 $a=3,c=5,e=5/3$，上焦点 $F_U(0,5)$，上准线 $y=9/5$。
已知点的纵坐标 $y=6$，它到上准线的垂距是
\[
\left|6-\frac95\right|=\frac{21}{5},
\qquad PF_U=\frac53\cdot\frac{21}{5}=\boxed7.
\]
此时 $x^2=16(36/9-1)=48$，
勾股回验 $PF_U^2=x^2+(6-5)^2=49$。

\medskip
\textbf{例 8（新）：由焦准比反求离心率与开口。}
已知某双曲线上的点 $P$ 到一组\textbf{对应}焦点、准线的距离
分别为 $12$ 和 $8$，则
\[
\boxed{e=12/8=3/2>1},\qquad
\boxed{\frac ba=\sqrt{e^2-1}=\frac{\sqrt5}{2}}.
\]
比值可以确定 $e$ 与两渐近线斜率的绝对值，但\textbf{不能仅凭这两段长度}
确定双曲线的实际大小与位置；还需要其他条件。
\end{examplebox}
